% Luc Destombe --- licence CC BY-NC-SA 4.0
% https://creativecommons.org/licenses/by-nc-sa/4.0/deed.fr
% Fichier autonome : compiler avec pdflatex, deux fois (signets, nombre de pages).
\def\lyceeversion{1.0}
\documentclass[a4paper,10pt]{article}
\makeatletter
\newif\iflycee@niveaux
\lycee@niveauxtrue
\newif\iflycee@cm
\newif\iflycee@palatino

\RequirePackage[utf8]{inputenc}
\RequirePackage[T1]{fontenc}
\RequirePackage[french]{babel}
\RequirePackage{etoolbox}

\newcommand{\ShorthandsOff}{\@ifundefined{shorthandoff}{}{\shorthandoff{;:!?}}}
\newcommand{\ShorthandsOn}{\@ifundefined{shorthandon}{}{\shorthandon{;:!?}}}
\ShorthandsOff
\AtBeginDocument{\ShorthandsOn}
\AtBeginEnvironment{tikzpicture}{\ShorthandsOff}

\RequirePackage{geometry}
\RequirePackage{fancyhdr}
\RequirePackage{lastpage}
\RequirePackage{multicol}
\RequirePackage{array}
\RequirePackage{enumitem}
\RequirePackage{xcolor}
\RequirePackage{needspace}

\RequirePackage{amsmath,amssymb}
\iflycee@cm\else\RequirePackage{fourier}\fi
\DeclareMathAlphabet{\mathcal}{OMS}{cmsy}{m}{n}
\renewcommand{\ttdefault}{lmtt}
\begingroup\catcode`\_=\active \catcode`\^=\active
\global\let_\sb \global\let^\sp
\endgroup
\AtBeginDocument{\catcode`\_=12 \mathcode`\_="8000
  \catcode`\^=12 \mathcode`\^="8000 }
\newlength{\retraitcalculs}
\setlength{\retraitcalculs}{2em}

\RequirePackage{tikz}
\usetikzlibrary{arrows.meta,calc,babel,shapes.misc}
\RequirePackage[most]{tcolorbox}
\tcbuselibrary{skins,breakable}

\definecolor{coulDef}{HTML}{1F4E79}
\definecolor{coulProp}{HTML}{7030A0}
\definecolor{coulRem}{HTML}{C55A11}
\definecolor{coulEx}{HTML}{2E7D32}
\definecolor{coulExo}{HTML}{B03A2E}
\definecolor{coulPart}{HTML}{1F4E79}
\definecolor{coulMeth}{HTML}{1B6E4F}
\definecolor{coulCourbe}{HTML}{0B5ED7}
\definecolor{coulCourbeB}{HTML}{C00000}
\definecolor{coulCourbeC}{HTML}{2E7D32}
\definecolor{coulGrille}{HTML}{8C8C8C}
\definecolor{coulRep}{HTML}{1B6E4F}
\definecolor{coulBar}{HTML}{2E7D32}

\newcommand{\R}{\mathbb{R}}
\newcommand{\Rs}{\mathbb{R}^{*}}
\newcommand{\N}{\mathbb{N}}
\newcommand{\Ns}{\mathbb{N}^{*}}
\newcommand{\Z}{\mathbb{Z}}
\newcommand{\Q}{\mathbb{Q}}
\renewcommand{\le}{\leqslant}
\renewcommand{\ge}{\geqslant}
\renewcommand{\approx}{\simeq}
\newcommand{\vect}[1]{\overrightarrow{#1}}
\newcommand{\norme}[1]{\left\lVert #1 \right\rVert}
\newcommand{\intervalle}[2]{\left[#1\,;#2\right]}
\RequirePackage{cancel}
\renewcommand{\CancelColor}{\color{coulExo}}
\newcommand{\fois}[1]{\times{\color{coulExo}#1}}
\newcommand{\barre}[1]{\cancel{#1}}
\newlength{\lycee@eqh}\newlength{\lycee@eqd}
\newsavebox{\lycee@eqa}\newsavebox{\lycee@eqb}\newsavebox{\lycee@eqc}
\newcommand{\lycee@eqcol}[1]{\hbox{$\renewcommand{\arraystretch}{1.35}%
  \begin{array}[t]{@{}>{\displaystyle}r@{}>{\displaystyle{}}l@{}}#1\end{array}$}}
\newcommand{\lycee@eqmesure}[1]{%
  \ifdim\ht#1>\lycee@eqh\setlength{\lycee@eqh}{\ht#1}\fi
  \ifdim\dp#1>\lycee@eqd\setlength{\lycee@eqd}{\dp#1}\fi}
\newcommand{\lycee@eqtrait}{\hspace{1.2em}ou\hspace{1.2em}}
\newcommand{\lycee@eqiff}{$\iff$\hspace{0.8em}}
\newcommand{\eqproduit}[2]{\par\smallskip\noindent C'est une équation produit :
  $A\times B=0 \iff A=0$ ou $B=0$.\par\nopagebreak\smallskip
  \begingroup\setlength{\lycee@eqh}{0pt}\setlength{\lycee@eqd}{0pt}%
  \setbox\lycee@eqa=\lycee@eqcol{#1}\setbox\lycee@eqb=\lycee@eqcol{#2}%
  \lycee@eqmesure\lycee@eqa\lycee@eqmesure\lycee@eqb%
  \noindent\hspace*{\retraitcalculs}\lycee@eqiff\usebox\lycee@eqa\lycee@eqtrait\usebox\lycee@eqb\endgroup\par\smallskip}
\newcommand{\eqproduittrois}[3]{\par\smallskip\noindent C'est une équation produit :
  $A\times B\times C=0 \iff A=0$ ou $B=0$ ou $C=0$.\par\nopagebreak\smallskip
  \begingroup\setlength{\lycee@eqh}{0pt}\setlength{\lycee@eqd}{0pt}%
  \setbox\lycee@eqa=\lycee@eqcol{#1}\setbox\lycee@eqb=\lycee@eqcol{#2}\setbox\lycee@eqc=\lycee@eqcol{#3}%
  \lycee@eqmesure\lycee@eqa\lycee@eqmesure\lycee@eqb\lycee@eqmesure\lycee@eqc%
  \noindent\hspace*{\retraitcalculs}\lycee@eqiff\usebox\lycee@eqa\lycee@eqtrait\usebox\lycee@eqb\lycee@eqtrait\usebox\lycee@eqc\endgroup
  \par\smallskip}

\tikzset{grille/.style={coulGrille,line width=0.4pt}}
\newcommand{\quadrillage}[4]{\draw[grille,step=1] (#1,#2) grid (#3,#4);}
\tikzset{
  croix/.style={cross out,draw,line width=0.6pt,inner sep=0pt,minimum size=4pt},
  croixdroite/.style={croix,rotate=45,line width=0.8pt},
}
\newcommand{\pt}[3]{\path (#1) node[croix]{} node[#2,font=\small,inner sep=2.5pt]{$#3$};}

\newcommand{\lycee@repere}[4]{%
  \draw[grille,step=1] (#1,#2) grid (#3,#4);
  \draw[-{Stealth[length=2mm]},thick] (#1,0)--({#3+0.4},0) node[below left=-1pt]{$x$};
  \draw[-{Stealth[length=2mm]},thick] (0,#2)--(0,{#4+0.4}) node[below left=-1pt]{$y$};
  \node[below left,font=\scriptsize,inner sep=1.5pt] at (0,0) {$0$};}
\newcommand{\lycee@gradx}[1]{\foreach \k/\l in {#1}{%
  \draw (\k,0.07)--(\k,-0.07) node[below,font=\scriptsize,inner sep=1.5pt]{$\l$};}}
\newcommand{\lycee@grady}[1]{\foreach \k/\l in {#1}{%
  \draw (0.07,\k)--(-0.07,\k) node[left,font=\scriptsize,inner sep=1.5pt]{$\l$};}}
\newcommand{\lycee@intff}[2]{\left[#1\,;#2\right]}
\newcommand{\lycee@intfo}[2]{\left[#1\,;#2\right[}
\newcommand{\lycee@intof}[2]{\left]#1\,;#2\right]}
\newcommand{\lycee@intoo}[2]{\left]#1\,;#2\right[}
\newcommand{\lycee@ens}[1]{\left\{#1\right\}}
\AtBeginDocument{%
  \providecommand{\repere}{\lycee@repere}%
  \providecommand{\gradx}{\lycee@gradx}%
  \providecommand{\grady}{\lycee@grady}%
  \providecommand{\Cf}{\mathcal{C}\sb{f}}%
  \providecommand{\Cg}{\mathcal{C}\sb{g}}%
  \providecommand{\intff}{\lycee@intff}%
  \providecommand{\intfo}{\lycee@intfo}%
  \providecommand{\intof}{\lycee@intof}%
  \providecommand{\intoo}{\lycee@intoo}%
  \providecommand{\ens}{\lycee@ens}}

\newlist{questions}{enumerate}{3}
\setlist[questions,1]{label=\arabic*\textdegree),leftmargin=*,itemsep=2pt,topsep=3pt}
\setlist[questions,2]{label=\alph*),leftmargin=*,itemsep=1pt,topsep=2pt}
\setlist[questions,3]{label=\roman*),leftmargin=*,itemsep=1pt,topsep=1pt}
\setlist[itemize]{label=\textbullet,leftmargin=*,itemsep=2pt,topsep=3pt}

\newcommand{\besoinplace}[1]{\par\penalty\@M\begingroup
  \dimen@=#1\relax\dimen@ii=\pagegoal\advance\dimen@ii-\pagetotal
  \ifdim\dimen@>\dimen@ii\ifdim\dimen@ii>\z@\vfil\fi\break\fi\endgroup}

\newcommand{\lycee@pied}{\thepage/\pageref{LastPage}}
\newcommand{\entete}[2]{%
  \pagestyle{fancy}\fancyhf{}%
  \fancyhead[L]{\footnotesize\color{coulPart}\textbf{#1}}%
  \fancyhead[R]{\footnotesize\color{coulPart}\textbf{#2}}%
  \fancyfoot[C]{\footnotesize\lycee@pied}%
  \renewcommand{\headrulewidth}{0.4pt}%
  \renewcommand{\headrule}{\hbox to\headwidth{%
    \color{coulPart!40}\leaders\hrule height \headrulewidth\hfill}}}

\RequirePackage{ccicons}
\newcommand{\lycee@licence}{{\scriptsize\color{gray}%
  \copyright\ Luc Destombe --- \ccbyncsa\ CC BY-NC-SA 4.0}}
\newif\iflycee@licence \lycee@licencetrue
\newcommand{\sanslicence}{\lycee@licencefalse}
\AtBeginDocument{\iflycee@licence\AddToHookNext{shipout/foreground}{%
  \put(0,0){\rlap{\kern\dimexpr1in+\hoffset+\oddsidemargin+\textwidth\relax
    \llap{\raisebox{-\dimexpr1in+\voffset+\topmargin+\headheight+\headsep
      +\textheight+\footskip\relax}{\lycee@licence}}}}}\fi}

\newcommand{\lycee@site}{faire-des-maths.fr}
\newcommand{\lycee@version}{{\scriptsize\color{gray}\lycee@site
  \ifdefined\lyceeversion\ --- Version~\lyceeversion\fi}}
\AtBeginDocument{\iflycee@licence\AddToHookNext{shipout/foreground}{%
  \put(0,0){\rlap{\kern\dimexpr1in+\hoffset+\oddsidemargin\relax
    \raisebox{-\dimexpr1in+\voffset+\topmargin+\headheight+\headsep
      +\textheight+\footskip\relax}{\lycee@version}}}}\fi}

\newcommand{\formulesserrees}{%
  \setlength{\abovedisplayskip}{3pt}\setlength{\belowdisplayskip}{3pt}%
  \setlength{\abovedisplayshortskip}{2pt}\setlength{\belowdisplayshortskip}{3pt}}

\setlength{\parindent}{0pt}

\RequirePackage[implicit=false,hidelinks,pdfencoding=auto,psdextra,bookmarksopen,
  bookmarksopenlevel=1,pdfcreator={},pdfproducer={}]{hyperref}
\RequirePackage{bookmark}
\hypersetup{pdfauthor={Luc Destombe},pdfkeywords={CC BY-NC-SA 4.0}}
\newcounter{lycee@signet}
\newcommand{\signet}[3][]{\stepcounter{lycee@signet}%
  \edef\lycee@dest{\if\relax\detokenize{#1}\relax
    signet.\arabic{lycee@signet}\else#1\fi}%
  \hypertarget{\lycee@dest}{}\bookmark[level=#2,dest=\lycee@dest]{#3}}
\pdfstringdefDisableCommands{%
  \def\R{R}\def\N{N}\def\Z{Z}\def\Q{Q}\def\vect#1{#1}\def\Cf{Cf}\def\Cg{Cg}\def\fois#1{×#1}%
  \def\barre#1{#1}\def\intervalle#1#2{[#1;#2]}\def\intff#1#2{[#1;#2]}%
  \def\textsuperscript#1{#1}\def\dfrac#1#2{#1/#2}\def\frac#1#2{#1/#2}%
  \def\vec#1{#1⃗}}%
\RequirePackage[official]{eurosym}

\geometry{top=0.8cm,bottom=1.3cm,left=1cm,right=1cm,footskip=0.6cm}

\pagestyle{fancy}
\fancyhf{}
\renewcommand{\headrulewidth}{0pt}
\fancyfoot[C]{\footnotesize Page \thepage{} / \pageref{LastPage}}

\newcommand{\titrefiche}[2]{%
  \begin{center}
    \Large\textbf{#1}\\[2pt]
    \large\textbf{#2}\\[2pt]
  \end{center}
  \vspace{0.10cm}}

\setlength{\columnseprule}{0.4pt}
\setlength{\columnsep}{15pt}
\renewcommand{\columnseprulecolor}{\color{black!45}}
\setlength{\multicolsep}{2pt}

\newcounter{partiefiche}
\renewcommand{\thepartiefiche}{\Roman{partiefiche}}
\newif\ifapresbandeau
\newcommand{\lycee@bandeau}[1]{%
  \par\penalty-600
  \vspace{0.08cm}%
  \noindent\colorbox{coulPart}{%
    \parbox{\dimexpr\linewidth-2\fboxsep\relax}{%
      \color{white}\bfseries\raggedright\strut#1\strut}}%
  \nopagebreak\par\nopagebreak
  \vspace{0.06cm}\nopagebreak
  \apresbandeautrue}
\newcommand{\partiefiche}[1]{%
  \refstepcounter{partiefiche}%
  \lycee@bandeau{\signet{1}{\thepartiefiche{} --- #1}\thepartiefiche{} --- #1}}
\newcommand{\partiefichesansnum}[1]{\lycee@bandeau{\signet{1}{#1}#1}}

\newcounter{exo}
\newlength{\espaceexo}\setlength{\espaceexo}{7pt}
\newlength{\espacetitre}\setlength{\espacetitre}{2pt}
\newcommand{\exo}[1][\relax]{%
  \par
  \ifapresbandeau\apresbandeaufalse\else\penalty-150\addvspace{\espaceexo}\fi
  \refstepcounter{exo}%
  \noindent\signet[exercice-\arabic{exo}]{2}{Exercice \arabic{exo}}%
  \textbf{\color{coulExo}\underline{Exercice \theexo}}%
  \ifx\relax#1\relax\else\textbf{ : #1}\fi
  \par\nopagebreak\vspace{\espacetitre}\nopagebreak
  \ignorespaces}

\setlist[enumerate]{nosep,leftmargin=*,itemsep=2pt,topsep=2pt}
\setlist[itemize]{nosep,leftmargin=*,topsep=2pt}
\newcommand{\choix}[3]{\par\hspace*{1em}%
  \makebox[0.3\linewidth][l]{a) #1}\makebox[0.3\linewidth][l]{b) #2}c) #3}
\newcommand{\deux}[2]{\par\hspace*{0.6em}%
  \makebox[0.48\linewidth][l]{#1}#2}
\newcommand{\trois}[3]{\par\hspace*{0.6em}%
  \makebox[0.32\linewidth][l]{#1}\makebox[0.32\linewidth][l]{#2}#3}

  \renewcommand{\exo}[1][\relax]{%
    \ifapresbandeau\apresbandeaufalse\else\par\penalty-150\fi
    \refstepcounter{exo}%
    \vspace{0.05cm}%
    \noindent\signet[exercice-\arabic{exo}]{2}{Exercice \arabic{exo}}%
    \textbf{\color{coulExo}\underline{Exercice \theexo}}%
    \ifx\relax#1\relax\else\textbf{ : #1}\fi%
    \par\nopagebreak
    \ignorespaces}
  \newcommand{\rappel}[1]{%
    \par\vspace{0.06cm}\nopagebreak
    \noindent\fcolorbox{black!35}{black!5}{%
      \parbox{\dimexpr\linewidth-2\fboxsep-2\fboxrule\relax}{%
        \small\textbf{Rappel.} #1}}%
    \par\vspace{0.06cm}\nopagebreak}
  \setlist[enumerate]{nosep,leftmargin=*}
  \setlist[itemize]{nosep,leftmargin=*}
  \setlength{\multicolsep}{3pt plus 1pt minus 1pt}
  \setlength{\premulticols}{10pt}
  \setlength{\postmulticols}{10pt}
  \newlist{expr}{enumerate}{1}
  \setlist[expr]{label={\Alph*\ $=$},nosep,leftmargin=*,itemsep=0pt}

\renewenvironment{center}{\par\addvspace{3pt}\centering}{\par\addvspace{3pt}}
\formulesserrees
\AtBeginDocument{\formulesserrees}
\clubpenalty=10000
\widowpenalty=10000
\displaywidowpenalty=10000
\brokenpenalty=10000
\setlength{\parskip}{0pt}

\RequirePackage{ragged2e}
\setlength{\RaggedRightRightskip}{0pt plus 1fil}
\AtBeginDocument{\RaggedRight\binoppenalty=10000 \relpenalty=10000 }
\g@addto@macro\@parboxrestore{\RaggedRight}
\makeatother

\newtcblisting{python}{%
  listing only,enhanced,
  colback=black!3,colframe=black!35,boxrule=0.5pt,arc=1pt,
  left=6pt,right=6pt,top=2pt,bottom=2pt,
  listing options={language=Python,basicstyle=\ttfamily\small,
    keywordstyle=\color{coulMeth}\bfseries,columns=fullflexible,
    keepspaces=true,showstringspaces=false,aboveskip=0pt,belowskip=0pt}}

\newcommand{\qcm}[4]{\par\hspace*{0.6em}\makebox[0.48\linewidth][l]{a) #1}b) #2
  \par\hspace*{0.6em}\makebox[0.48\linewidth][l]{c) #3}d) #4}
\newcommand{\qcmligne}[4]{\par\hspace*{0.6em}\makebox[0.24\linewidth][l]{a) #1}%
  \makebox[0.24\linewidth][l]{b) #2}\makebox[0.24\linewidth][l]{c) #3}d) #4}

\begin{document}

\titrefiche{1\textsuperscript{re} STMG --- Variables aléatoires, loi de Bernoulli}{Fiche d'exercices du chapitre}

\begin{multicols}{2}\raggedcolumns

\partiefiche{Variable aléatoire}

\exo[Calculer des probabilités avec une variable aléatoire]
Au stand du club photo, on lance un dé équilibré à six faces : on gagne $3$~€ avec un $5$
ou un $6$, rien avec un $4$, et on perd $2$~€ sinon. On note $X$ le gain en euros.
\begin{enumerate}[label=\arabic*)]
  \item Donner la valeur de $X$ pour chaque issue, puis les valeurs prises par $X$.
  \item Calculer $P(X=3)$ et $P(X=-2)$.
  \item Traduire par une phrase l'événement $\{X\ge 0\}$, puis calculer sa probabilité.
  \item Calculer $P(X\le 0)$.
\end{enumerate}

\exo[En phrase complète]
Un commercial rencontre $10$ clients dans la journée. On note $X$ le nombre de clients
qui passent une commande.
\begin{enumerate}[label=\arabic*)]
  \item Écrire avec $X$ chacun des événements suivants.
        \begin{enumerate}[label=\alph*)]
          \item Exactement $7$ clients passent une commande.
          \item Au plus $3$ clients passent une commande.
          \item Au moins $8$ clients passent une commande.
          \item Moins de $2$ clients passent une commande.
          \item Plus de $6$ clients passent une commande.
          \item Tous les clients passent une commande.
        \end{enumerate}
  \item Traduire par une phrase les événements $\{X=0\}$, $\{X\le 9\}$ et $\{X\ge 1\}$.
  \item Écrire la probabilité de l'événement \og entre $4$ et $6$ clients passent une commande,
        bornes comprises\fg{} comme une somme de probabilités $P(X=\dots)$.
\end{enumerate}

\exo[Des journées promotion]
Pour chacun des $3$ prochains jours, une boutique en ligne tire au sort si le jour sera
une journée promotion ($P$) ou non ($N$), avec la même probabilité. Les $8$ issues
($PPP$ ; $PPN$ ; \dots) ont la même probabilité. On note $X$ le nombre de journées
promotion.
\begin{enumerate}[label=\arabic*)]
  \item Écrire les $8$ issues dans un tableau, avec la valeur de $X$ pour chacune.
  \item Quelles valeurs $X$ peut-il prendre ?
  \item Calculer $P(X=2)$ et $P(X=0)$.
  \item Traduire par une phrase l'événement $\{X\ge 1\}$, puis calculer sa probabilité.
\end{enumerate}

\exo[À l'ouverture d'une boutique]
Une urne contient $10$ boules numérotées de $1$ à $10$. Une partie coûte $1$~€ : on tire
une boule au hasard ; on reçoit $5$~€ si son numéro est un multiple de $5$, $1$~€ si c'est
un multiple de $3$, et rien sinon. On note $X$ le gain en euros, prix de la partie déduit.
\begin{enumerate}[label=\arabic*)]
  \item Donner la valeur de $X$ pour chaque numéro, puis les valeurs prises par $X$.
  \item Calculer $P(X=4)$ et $P(X=-1)$.
  \item Traduire par une phrase l'événement $\{X\ge 0\}$, puis calculer sa probabilité.
  \item Calculer $P(X<4)$ à l'aide de l'événement contraire.
\end{enumerate}

\partiefiche{Loi de probabilité d'une variable aléatoire}

\exo[Déterminer et utiliser une loi]
\begin{enumerate}[label=\alph*)]
  \item Pour financer son festival, une association vend $200$ tickets de tombola, à
        $3$~€ chacun. $2$ tickets font gagner un lot de $50$~€, $10$ tickets un lot de $10$~€, les
        autres ne font rien gagner. On achète un ticket au hasard ; on note $X$ le gain en
        euros, prix du ticket déduit. Déterminer la loi de probabilité de $X$.
  \item On lance deux dés équilibrés à quatre faces, numérotées de $1$ à $4$. On note $Y$
        le plus petit des deux nombres obtenus (s'ils sont égaux, ce nombre). Déterminer la
        loi de probabilité de $Y$.
  \item Dans une boulangerie, on note $Z$ le nombre de clients qui attendent devant la
        caisse, à un instant pris au hasard. Sa loi de probabilité
        est la suivante. Calculer $a$, puis $P(Z\le 1)$ et la probabilité qu'au moins $2$
        clients attendent.
        \begin{center}
        \renewcommand{\arraystretch}{1.4}
        \begin{tabular}{|c|c|c|c|c|c|}
        \hline
        $z_i$ & $0$ & $1$ & $2$ & $3$ & $4$\\
        \hline
        $P(Z=z_i)$ & $0{,}35$ & $0{,}30$ & $a$ & $0{,}10$ & $0{,}05$\\
        \hline
        \end{tabular}
        \end{center}
\end{enumerate}

\exo[Des réclamations]
Un service après-vente reçoit chaque jour entre $0$ et $5$ réclamations. On note $X$ le
nombre de réclamations un jour pris au hasard ; sa loi de probabilité est la
suivante.
\begin{center}
\renewcommand{\arraystretch}{1.4}
\begin{tabular}{|c|c|c|c|c|c|c|}
\hline
$x_i$ & $0$ & $1$ & $2$ & $3$ & $4$ & $5$\\
\hline
$P(X=x_i)$ & $0{,}05$ & $0{,}15$ & $0{,}30$ & $a$ & $0{,}15$ & $0{,}10$\\
\hline
\end{tabular}
\end{center}
\begin{enumerate}[label=\arabic*)]
  \item Calculer $a$.
  \item Traduire par une phrase les événements $\{X\le 2\}$ et $\{X\ge 4\}$, puis
        calculer leurs probabilités.
  \item Calculer la probabilité que le service reçoive plus de $2$ réclamations.
  \item Pour un autre service, un stagiaire écrit le tableau ci-dessous. Expliquer
        pourquoi ce ne peut pas être une loi de probabilité.
        \begin{center}
        \renewcommand{\arraystretch}{1.4}
        \begin{tabular}{|c|c|c|c|c|}
        \hline
        $y_i$ & $0$ & $1$ & $2$ & $3$\\
        \hline
        $P(Y=y_i)$ & $0{,}2$ & $0{,}3$ & $0{,}3$ & $0{,}3$\\
        \hline
        \end{tabular}
        \end{center}
\end{enumerate}

\exo[Un dé publicitaire]
Un dé publicitaire à six faces, numérotées de $1$ à $6$, n'est pas équilibré :
chaque face paire a deux fois plus de chances de sortir que chaque face impaire. On note
$p$ la probabilité d'obtenir $1$, qui est aussi celle d'obtenir $3$ et celle d'obtenir
$5$, et $X$ le numéro obtenu.
\begin{enumerate}[label=\arabic*)]
  \item Exprimer en fonction de $p$ la probabilité d'obtenir $2$.
  \item Montrer que $9p=1$, puis calculer $p$.
  \item Déterminer la loi de probabilité de $X$.
  \item Calculer $P(X\ge 4)$, puis la probabilité d'obtenir un nombre pair.
\end{enumerate}

\partiefiche{Espérance}

\exo[Calculer et interpréter une espérance]
\begin{enumerate}[label=\alph*)]
  \item Un drone livre des colis dans une zone rurale. On note $X$ le nombre de colis
        livrés en une heure prise au hasard ; sa loi est la suivante. Calculer $E(X)$ et
        interpréter le résultat.
        \begin{center}
        \renewcommand{\arraystretch}{1.4}
        \begin{tabular}{|c|c|c|c|c|}
        \hline
        $x_i$ & $0$ & $1$ & $2$ & $3$\\
        \hline
        $P(X=x_i)$ & $0{,}10$ & $0{,}30$ & $0{,}40$ & $0{,}20$\\
        \hline
        \end{tabular}
        \end{center}
  \item À une tombola, on note $Y$ le gain en euros d'un visiteur, prix du ticket
        déduit ; sa loi est la suivante. Calculer $E(Y)$ et interpréter le résultat, pour
        le visiteur puis pour l'association qui organise la tombola.
        \begin{center}
        \renewcommand{\arraystretch}{1.4}
        \begin{tabular}{|c|c|c|c|}
        \hline
        $y_i$ & $-2$ & $3$ & $18$\\
        \hline
        $P(Y=y_i)$ & $0{,}85$ & $0{,}10$ & $0{,}05$\\
        \hline
        \end{tabular}
        \end{center}
\end{enumerate}

\exo[Deux stands]
À la fête du lycée, deux stands proposent un jeu. On note $X$ le gain en euros au stand A
et $Y$ le gain en euros au stand B ; leurs lois sont les suivantes.
\begin{center}
\renewcommand{\arraystretch}{1.4}
\begin{tabular}{|c|c|c|c|}
\hline
$x_i$ & $-1$ & $2$ & $5$\\
\hline
$P(X=x_i)$ & $0{,}75$ & $0{,}20$ & $0{,}05$\\
\hline
\end{tabular}
\qquad
\begin{tabular}{|c|c|c|c|}
\hline
$y_i$ & $-2$ & $4$ & $20$\\
\hline
$P(Y=y_i)$ & $0{,}85$ & $0{,}10$ & $0{,}05$\\
\hline
\end{tabular}
\end{center}
\begin{enumerate}[label=\arabic*)]
  \item Calculer $E(X)$ et $E(Y)$.
  \item Pour quelqu'un qui joue souvent, quel stand est le moins défavorable ? Justifier.
  \item Calculer $P(X>0)$ et $P(Y>0)$.
  \item Un ami affirme : \og Au stand B, on peut gagner $20$~€ : c'est le meilleur.\fg{}
        Qu'en penser ? Répondre par une phrase.
\end{enumerate}

\exo[Au magasin de cycles]
Dans un magasin de cycles, $20\,\%$ des clients achètent un vélo électrique, à
$1\,200$~€. Parmi eux, $60\,\%$ achètent aussi un casque, à $50$~€ ; parmi les autres
clients, $10\,\%$ achètent un casque seul. On choisit un client au hasard ; on note $V$
\og le client achète un vélo\fg{}, $C$ \og le client achète un casque\fg{} et $X$ la somme
qu'il dépense, en euros, pour les vélos et les casques.
\begin{enumerate}[label=\arabic*)]
  \item Construire l'arbre pondéré de la situation.
  \item Quelles valeurs $X$ peut-il prendre ?
  \item Déterminer la loi de probabilité de $X$.
  \item Calculer $E(X)$ et interpréter le résultat.
  \item Le magasin reçoit environ $400$ clients par mois. Quel chiffre d'affaires mensuel
        peut-il prévoir pour les vélos et les casques ?
\end{enumerate}

\exo[Le coût d'une trottinette]
Une entreprise fabrique des trottinettes électriques ; chacune lui coûte $200$~€. Un
contrôle repère un défaut de frein (réparation : $20$~€), un défaut de batterie
(réparation : $30$~€) ou les deux (réparation : $50$~€). Sur l'ensemble de la production,
$92\,\%$ des trottinettes n'ont aucun défaut, $4\,\%$ ont seulement le défaut de frein,
$3\,\%$ seulement le défaut de batterie et $1\,\%$ les deux. On choisit une trottinette
au hasard et on note $X$ son coût total en euros, réparation comprise.
\begin{enumerate}[label=\arabic*)]
  \item Quelles valeurs $X$ peut-il prendre ?
  \item Déterminer la loi de probabilité de $X$.
  \item Calculer $E(X)$ et interpréter le résultat pour l'entreprise.
  \item Si l'entreprise vend chaque trottinette $200$~€, gagne-t-elle de l'argent en
        moyenne ?
  \item À quel prix doit-elle la vendre pour gagner en moyenne $50$~€ par trottinette ?
\end{enumerate}

\partiefiche{Loi de Bernoulli}

\exo[Reconnaître une loi de Bernoulli]
Pour chaque variable aléatoire $X$, $Y$ et $Z$, dire si elle suit une loi de Bernoulli.
Si oui, donner son paramètre, sa loi et son espérance.
\begin{enumerate}[label=\alph*)]
  \item Un bus électrique arrive à l'heure à son arrêt avec la probabilité $0{,}85$. On
        observe un de ses passages ; $X$ vaut $1$ s'il arrive à l'heure, $0$ sinon.
  \item On lance un dé équilibré à douze faces, numérotées de $1$ à $12$ ; $Y$ vaut $1$ si
        on obtient un multiple de $3$, $0$ sinon.
  \item On lance le même dé ; $Z$ est le numéro obtenu.
\end{enumerate}

\exo[Épreuve de Bernoulli ou non ?]
Pour chaque expérience, dire par une phrase s'il s'agit d'une épreuve de Bernoulli. Si
oui, préciser le succès choisi et, quand on peut la calculer, sa probabilité.
\begin{enumerate}[label=\alph*)]
  \item Un magasin vend des vélos de taille S, M ou L ; on choisit au hasard un vélo
        vendu et on note sa taille.
  \item On lance une pièce équilibrée et on regarde si elle tombe sur Pile.
  \item On choisit au hasard un jour de l'année et on note son mois.
  \item On choisit au hasard un colis expédié par une boutique en ligne et on regarde
        s'il arrive en bon état.
  \item On tire une carte au hasard dans un jeu de $32$ cartes et on regarde si c'est un
        as.
\end{enumerate}

\exo[Le défi du dodécaèdre]
Pour une partie, on paie $2$~€ et on lance un dé équilibré à douze faces, numérotées de
$1$ à $12$. Si le numéro est un multiple de $4$, on reçoit $7$~€ ; sinon, rien. On note
$S$ l'événement \og le numéro est un multiple de $4$\fg{}.
\begin{enumerate}[label=\arabic*)]
  \item Calculer $P(S)$.
  \item On note $X$ la variable aléatoire qui vaut $1$ si $S$ est réalisé, $0$ sinon.
        Donner la loi de probabilité de $X$ et son espérance.
  \item On note $G$ le gain du joueur en euros, prix de la partie déduit. Déterminer la
        loi de probabilité de $G$. La variable $G$ suit-elle une loi de Bernoulli ?
  \item Calculer $E(G)$ et interpréter le résultat.
\end{enumerate}

\partiefiche{Répétition d'épreuves de Bernoulli}

\exo[Calculer une probabilité avec un arbre]
\begin{enumerate}[label=\alph*)]
  \item Sur une plateforme de vente, une commande est livrée à l'heure avec la
        probabilité $0{,}8$, indépendamment des autres. On suit $2$ commandes ; on note
        $X$ le nombre de commandes livrées à l'heure. Construire l'arbre, puis calculer $P(X=2)$ et $P(X=1)$.
  \item Un livreur remet un colis du premier coup avec la probabilité $0{,}9$,
        indépendamment des autres colis. Il livre $3$ colis ; on note $Y$ le nombre de
        colis remis du premier coup. Construire l'arbre, puis calculer $P(Y=3)$, $P(Y=2)$ et
        $P(Y\ge 1)$.
\end{enumerate}

\exo[Trois lancers d'une pièce]
On lance trois fois une pièce équilibrée ; on note $X$ le nombre de Pile obtenus.
\begin{enumerate}[label=\arabic*)]
  \item Construire l'arbre de l'expérience, avec les probabilités sur les branches.
  \item Calculer $P(X=3)$ et $P(X=2)$.
  \item Écrire avec $X$ l'événement \og on obtient au moins une fois Face\fg{}, puis
        calculer sa probabilité.
  \item Déterminer la loi de probabilité de $X$, puis son espérance.
\end{enumerate}

\exo[Une agence de location]
Dans une agence de location, chaque voiture est en entretien avec la probabilité
$0{,}2$, indépendamment des autres. Un jour pris au hasard, on observe $3$ voitures ; on
note $X$ le nombre de voitures en entretien.
\begin{enumerate}[label=\arabic*)]
  \item Construire l'arbre de l'expérience.
  \item Déterminer la loi de probabilité de $X$.
  \item Calculer $E(X)$, puis comparer avec $3\times 0{,}2$.
  \item Calculer la probabilité qu'au plus une voiture soit en entretien.
\end{enumerate}

\exo[Fléchettes électroniques]
À chaque lancer, une fléchette atteint la zone rouge avec la probabilité $0{,}6$,
indépendamment des autres lancers : elle rapporte alors $2$ points ; sinon, elle en fait
perdre $1$. Un joueur lance $2$ fléchettes ; on note $X$ son score.
\begin{enumerate}[label=\arabic*)]
  \item Construire l'arbre de l'expérience.
  \item Quelles valeurs $X$ peut-il prendre ?
  \item Déterminer la loi de probabilité de $X$.
  \item Calculer $P(X>0)$ et interpréter le résultat.
  \item Calculer $E(X)$ et interpréter le résultat.
\end{enumerate}

\partiefiche{Simulation d'échantillons}

\exo[Simuler à la main et à l'ordinateur]
Un jour pris au hasard, une boutique en ligne réalise plus de $50$ ventes avec la
probabilité $0{,}8$, indépendamment des autres jours.
\begin{enumerate}[label=\arabic*)]
  \item Compléter la fonction \texttt{jour()}, qui renvoie $1$ si la boutique réalise plus
        de $50$ ventes et $0$ sinon.
\begin{python}
from random import random

def jour():
    if random() < .....:
        return 1
    else:
        return .....
\end{python}
  \item Avec les chiffres au hasard ci-dessous (de $0$ à $7$ pour plus de $50$ ventes,
        $8$ ou $9$ sinon), simuler deux échantillons de $10$ jours, puis donner la
        fréquence des jours à plus de $50$ ventes dans chacun.
        \par\smallskip
        Échantillon 1 : \texttt{3 8 0 6 9 1 5 2 7 4} \qquad Échantillon 2 :
        \texttt{9 1 8 4 0 9 6 3 2 5}
  \item L'ordinateur a simulé $10$ échantillons de $20$ jours. Nombres de jours à plus de
        $50$ ventes : $16$ ; $15$ ; $17$ ; $16$ ; $14$ ; $18$ ; $16$ ; $15$ ; $17$ ; $16$.
        Donner la fréquence du premier échantillon. Pourquoi les fréquences sont-elles
        différentes ?
  \item Calculer la moyenne des $10$ fréquences, et la comparer à $0{,}8$.
\end{enumerate}

\exo[Lire un programme]
Dans une imprimerie, $5\,\%$ des livres présentent un défaut d'impression. On choisit un
livre au hasard ; on note $D$ l'événement \og le livre présente un défaut\fg{}.
\begin{python}
from random import random

def livre():
    if random() < 0.05:
        return 1
    else:
        return 0
\end{python}
\begin{enumerate}[label=\arabic*)]
  \item Donner $P(D)$.
  \item Expliquer ce que renvoie la fonction \texttt{livre()}, et avec quelle probabilité
        elle renvoie $1$.
  \item L'ordinateur a appelé $10\,000$ fois la fonction \texttt{livre()} : $493$ livres
        simulés ont un défaut. Calculer la fréquence observée des livres avec un défaut.
  \item Ce résultat est-il étonnant ? Expliquer.
\end{enumerate}

\exo[La moyenne des fréquences]
Dans une agence de voyages, un client choisit l'option \og bagage en soute\fg{} avec la
probabilité $0{,}3$. L'ordinateur a simulé $5$ échantillons de $100$ clients ; voici la
fréquence des clients qui choisissent l'option dans chacun :
\[
0{,}28 \qquad 0{,}33 \qquad 0{,}31 \qquad 0{,}27 \qquad 0{,}31
\]
\begin{enumerate}[label=\arabic*)]
  \item Calculer la moyenne de ces $5$ fréquences.
  \item On choisit un client au hasard ; $X$ vaut $1$ s'il choisit l'option, $0$ sinon.
        Quelle loi $X$ suit-elle ? Donner $E(X)$, puis comparer avec la question 1.
  \item Sur $500$ échantillons simulés, la moyenne des fréquences vaut $0{,}2997$.
        Interpréter ce résultat.
\end{enumerate}

\partiefiche{Fluctuation d'échantillonnage}

\exo[Fréquente ou rare ?]
Une pièce équilibrée tombe sur Pile avec la probabilité $p=0{,}5$. Sur $200$ échantillons
simulés de $100$ lancers, l'écart-type des fréquences de Pile vaut $s\approx 0{,}05$.
\begin{enumerate}[label=\arabic*)]
  \item Calculer $p-2s$ et $p+2s$.
  \item Léo lance $100$ fois sa pièce et obtient $57$ Pile. Est-ce une situation fréquente
        ou rare dans le modèle ?
  \item Une autre pièce donne $68$ Pile sur $100$ lancers. Même question ; que
        peut-on en penser ?
  \item Avec des échantillons de $400$ lancers, quelle valeur de $s$ peut-on prévoir ?
\end{enumerate}

\exo[Un générateur de chiffres]
Un générateur de chiffres $0$ ou $1$ doit produire des $1$ avec la probabilité $p=0{,}5$. On a
simulé $200$ paquets de $25$ chiffres ; l'histogramme donne le nombre de paquets pour chaque
nombre de chiffres égaux à $1$. Sur cette simulation, l'écart-type des fréquences des $1$ vaut
$s\approx 0{,}1$.
\begin{center}
\begin{tikzpicture}[x=0.4cm,y=0.075cm]
  \foreach \k/\n in {6/1,7/1,8/11,9/8,10/16,11/35,12/29,13/33,14/23,15/18,16/14,
                     17/3,18/6,19/2}{
    \draw[fill=coulCourbe!30] (\k-0.4,0) rectangle (\k+0.4,\n);
    \node[above,font=\scriptsize,inner sep=1pt] at (\k,\n) {$\n$};
    \node[below,font=\scriptsize,inner sep=2pt] at (\k,0) {$\k$};}
  \draw[-{Stealth[length=2mm]},thick] (5,0) -- (20.4,0);
  \draw[-{Stealth[length=2mm]},thick] (5,0) -- (5,41);
  \node[below left,font=\footnotesize] at (20.4,-4.5) {nombre de $1$ dans le paquet};
  \node[right,font=\footnotesize] at (5,40) {nombre de paquets};
\end{tikzpicture}
\end{center}
\begin{enumerate}[label=\arabic*)]
  \item Combien de paquets contiennent $12$ chiffres égaux à $1$ ? Quelle est la fréquence des
        $1$ dans ces paquets ?
  \item Calculer $p-2s$ et $p+2s$. Pour quels nombres de $1$ la fréquence d'un paquet
        est-elle dans $[p-2s\,;p+2s]$ ?
  \item Combien de paquets ont une fréquence hors de cet intervalle ? Quel pourcentage des
        $200$ paquets cela représente-t-il ? Est-ce conforme au cours ?
  \item Un paquet produit par le générateur contient $21$ chiffres égaux à $1$. Est-ce une
        situation fréquente ou rare dans le modèle ?
\end{enumerate}

\exo[La taille de l'échantillon]
Pour estimer la part des usagers satisfaits d'un réseau de tramway, un institut interroge
des usagers au hasard. On a simulé des échantillons de taille $n$ d'une loi de Bernoulli
de paramètre $p=0{,}5$ ; le tableau donne l'écart-type $s$ des fréquences observées.
\begin{center}
\renewcommand{\arraystretch}{1.4}
\begin{tabular}{|l|c|c|c|c|}
\hline
Taille $n$ & $25$ & $100$ & $400$ & $2\,500$\\
\hline
Écart-type $s$ & $0{,}10$ & $0{,}05$ & \dots & \dots\\
\hline
\end{tabular}
\end{center}
\begin{enumerate}[label=\arabic*)]
  \item Comment évolue $s$ quand $n$ est multiplié par $4$ ? Compléter la case de
        $n=400$.
  \item De $n=25$ à $n=2\,500$, par combien la taille est-elle multipliée ? Compléter la
        dernière case.
  \item Expliquer par une phrase pourquoi une enquête auprès de $2\,500$ usagers est plus
        précise qu'une enquête auprès de $100$ usagers.
\end{enumerate}

\exo[Des batteries longue durée]
Un fabricant annonce que $90\,\%$ de ses batteries de vélo supportent plus de $1\,000$
recharges. Sur $200$ échantillons simulés de $100$ batteries, avec $p=0{,}9$,
l'écart-type des fréquences vaut $s\approx 0{,}03$.
\begin{enumerate}[label=\arabic*)]
  \item On choisit une batterie au hasard ; $X$ vaut $1$ si elle supporte plus de
        $1\,000$ recharges, $0$ sinon. Quelle loi $X$ suit-elle ? Donner $E(X)$.
  \item Calculer $p-2s$ et $p+2s$.
  \item Le laboratoire A teste $100$ batteries : $87$ supportent plus de $1\,000$
        recharges. Est-ce une situation fréquente ou rare dans le modèle ?
  \item Le laboratoire B teste $100$ batteries d'un autre lot : $78$ seulement supportent
        plus de $1\,000$ recharges. Même question ; que peut-on penser de ce lot ?
  \item Le laboratoire B peut-il être certain que ce lot ne respecte pas l'annonce ?
        Répondre par une phrase.
\end{enumerate}

\partiefichesansnum{Exercices type épreuve anticipée}

\exo[QCM]
Pour chaque question, une seule réponse est exacte.
\begin{enumerate}[label=\arabic*.]
  \item Une variable aléatoire $X$ prend les valeurs $0$ ; $1$ et $2$, avec
        $P(X=0)=0{,}2$ et $P(X=1)=0{,}5$. Alors $P(X=2)$ vaut :
        \qcmligne{$0{,}7$}{$0{,}3$}{$0{,}5$}{$1$}
  \item Un dé est truqué : la face $5$ sort avec la probabilité $0{,}2$ et la face $6$ avec
        la probabilité $0{,}35$. On note $X$ le numéro obtenu. Alors $P(X\ge 5)$ vaut :
        \qcmligne{$\dfrac{2}{6}$}{$0{,}35$}{$0{,}55$}{$0{,}07$}
  \item La loi d'une variable aléatoire $X$ est : $P(X=0)=0{,}5$ ; $P(X=2)=0{,}3$ ;
        $P(X=10)=0{,}2$. L'espérance de $X$ vaut :
        \qcmligne{$2{,}6$}{$4$}{$12$}{$0{,}6$}
  \item Une variable aléatoire $X$ suit la loi de Bernoulli de paramètre $0{,}3$. Alors
        $E(X)$ vaut :
        \qcmligne{$0{,}7$}{$1$}{$0{,}3$}{$0{,}5$}
  \item Dans un magasin, $80\,\%$ des clients paient par carte. On choisit au hasard
        deux clients, de façon indépendante. La probabilité que les deux paient par carte
        est :
        \qcmligne{$1{,}6$}{$0{,}64$}{$0{,}8$}{$0{,}16$}
  \item Pour un modèle de paramètre $p=0{,}5$, l'écart-type des fréquences vaut
        $s\approx 0{,}05$. Une fréquence observée de $0{,}62$ est :
        \par\hspace*{0.6em}a) une situation fréquente dans le modèle
        \par\hspace*{0.6em}b) une situation rare dans le modèle
        \par\hspace*{0.6em}c) égale à $p$
        \par\hspace*{0.6em}d) impossible dans le modèle
\end{enumerate}

\exo[Vrai ou faux ?]
Dans une boutique de souvenirs, on note $X$ le nombre d'articles achetés par un client
pris au hasard. Sa loi de probabilité est la suivante.
\begin{center}
\renewcommand{\arraystretch}{1.4}
\begin{tabular}{|c|c|c|c|c|}
\hline
$x_i$ & $1$ & $2$ & $3$ & $4$\\
\hline
$P(X=x_i)$ & $0{,}4$ & $0{,}3$ & $0{,}2$ & $0{,}1$\\
\hline
\end{tabular}
\end{center}
Pour chacune des quatre affirmations, dire si elle est vraie ou fausse, en justifiant.
\begin{enumerate}[label=\arabic*.]
  \item \emph{$P(X\ge 2)=0{,}3$.}
  \item \emph{En moyenne, un client achète $2$ articles.}
  \item \emph{Plus de la moitié des clients achètent au plus $2$ articles.}
  \item \emph{On choisit au hasard deux clients, de façon indépendante. La probabilité
        que chacun achète un seul article est égale à $0{,}16$.}
\end{enumerate}

\exo[Les nuits d'un hôtel]
Dans un hôtel, $30\,\%$ des clients restent $1$ nuit, $35\,\%$ restent $2$ nuits,
$20\,\%$ restent $3$ nuits, $10\,\%$ restent $4$ nuits et $5\,\%$ restent $5$ nuits. On
choisit un client au hasard et on note $X$ le nombre de nuits qu'il passe à l'hôtel.
\begin{enumerate}[label=\arabic*.]
  \item Déterminer la loi de probabilité de $X$.
  \item Calculer $P(X\ge 2)$. Interpréter ce résultat dans le contexte de l'exercice.
  \item Montrer que $E(X)=2{,}25$. Interpréter ce résultat.
  \item Une nuit est facturée $80$~€. Quelle recette moyenne l'hôtel peut-il attendre par
        client ?
  \item On choisit au hasard deux clients, de façon indépendante. Quelle est la
        probabilité que les deux ne restent qu'une nuit ? On pourra s'aider d'un arbre.
\end{enumerate}

\end{multicols}

\end{document}
